\section{背景}

高校文化祭や各種イベントにおいて、人気の高いアトラクションや展示への来場者集中は避けて通れない課題である。特に、処理能力（キャパシティ）に制限のあるインタラクティブな体験型コンテンツにおいては、待機列の極度な延伸が周辺通路の通行阻害や来場者の不満蓄積、さらには安全上のリスクを引き起こす。
これらの問題を解決するための有効な手法として、待機列の状況をリアルタイムに把握し、正確な待ち時間を試算・提示することで来場者の行動を分散させる方法が挙げられる。近年では、こうした案内業務を省力化・自動化するため、固定カメラの映像にコンピュータビジョン技術を適用し、待機列の混雑度を計測する取り組みが普及しつつある。
しかしながら、実際のイベント運用における待ち時間の算定や案内枠の管理においては、単に「エリア内に何人存在するか」という総人数の計測だけでは不十分である。来場者の多くは家族や友人、カップルといった「グループ（連れ）」単位で行動し、アトラクションもグループ単位で体験することが一般的であるため、システム側でも「人物個体」ではなく「グループ単位」で待機列の構造を把握することが強く求められている。